\section{Support of three-region operator factors}

Let $V$ be a finite set of sites, with local space $\C^q$ at each site.
For $S\subseteq V$, write $\mathcal H_S=(\C^q)^{\otimes S}$.
If $T,E\subseteq V$ are disjoint and $R=V\setminus(T\cup E)$,
restriction of configurations identifies the full space with
$\mathcal H_T\otimes\mathcal H_E\otimes\mathcal H_R$.
Write $P^g$ for an operator in these coordinates. Support on a set $D$
means that the operator acts trivially on the sites outside $D$.
The statements below give supporting locality algebra for
\cite[Proposition~5.7]{Arad2016RigorousRG}.

\begin{theorem}[Support of partial matrix elements]
    \label{thm:three_region_partial_support}
    \lean{TNLean.MPS.AreaLaw.operatorBlock_isSupportedOn_of_threeSplit_left,
        TNLean.MPS.AreaLaw.leftBlock_isSupportedOn_of_threeSplit_right}
    \leanok
    Suppose that $P$ is supported on $D\subseteq V$.
    If $P^g=U\otimes \Id_R$, then for any configurations $a,b$ on $T$,
    the matrix on $E$ with entries
    \begin{align}
        U^{\mathrm R}_{ab}(x,y)=U((a,x),(b,y))
        \label{eq:three_region_left_partial}
    \end{align}
    is supported on $D\cap E$, viewed as a subset of $E$.
    If instead $P^g=\Id_T\otimes V_0$, then for any configurations $r,s$
    on $R$, the matrix
    \begin{align}
        (V_0)^{\mathrm L}_{rs}(x,y)=V_0((x,r),(y,s))
        \label{eq:three_region_right_partial}
    \end{align}
    has the same support conclusion. The two fixed configurations in either
    statement need not agree outside $D$.
\end{theorem}
\begin{proof}\leanok
    For $q>0$, choose one fixed configuration on the unused region.
    The displayed partial matrix element is then an entry of $P$ with that
    configuration on both sides. A disagreement at a site of $E\setminus D$
    forces the entry to vanish. When the configurations agree there, the
    dependence on the remaining coordinates follows from the support of $P$.
    For $q=0$, the configuration set of $E$ is either empty or a singleton;
    the support assertion follows directly in both cases.
\end{proof}

\begin{theorem}[Equality from a middle identity lift]
    \label{thm:three_region_middle_lift_equality}
    \lean{TNLean.MPS.AreaLaw.eq_of_threeSplit_middle_lift_eq}
    \leanok
    Suppose that $E$ contains a specified site. If $W,M$ are operators on
    $\mathcal H_E$ and
    \begin{align}
        \Id_T\otimes W\otimes \Id_R=\Id_T\otimes M\otimes \Id_R,
        \label{eq:three_region_middle_lift_equality}
    \end{align}
    then $W=M$. No positive-local-dimension assumption is required.
\end{theorem}
\begin{proof}\leanok
    For $q>0$, evaluate the identity at equal fixed configurations on $T$ and
    $R$. For $q=0$, a configuration on $E$ would assign an element of the empty
    local alphabet to the specified site. Thus both matrices have empty index
    sets and are equal.
\end{proof}

\begin{theorem}[Canonical placement in the middle region]
    \label{thm:three_region_middle_embedding}
    \lean{TNLean.MPS.AreaLaw.embedOp_subtype_val_submatrix_threeSplit}
    \leanok
    For every operator $W$ on $\mathcal H_E$, its canonical extension to $V$
    has three-region form
    \begin{align}
        (\iota_E(W))^g=\Id_T\otimes W\otimes \Id_R.
        \label{eq:three_region_middle_embedding}
    \end{align}
    This includes empty regions and zero local dimension.
\end{theorem}
\begin{proof}\leanok
    The extension has entry $W(x,y)$ when the two configurations agree off
    $E$, and zero otherwise. In three-region coordinates, agreement off $E$
    is precisely equality of both the $T$ and $R$ configurations. These two
    equalities give the entries of the two identity matrices.
\end{proof}

\begin{theorem}[Two-window factors with their collars]
    \label{thm:three_region_two_window_collars}
    \lean{TNLean.MPS.AreaLaw.exists_threeSplit_two_window_factors_with_collars}
    \leanok
    \uses{def:chain_model,lem:peps_frame_split_forms,
        thm:three_region_partial_support}
    Consider a chain of $L$ sites and nonnegative integers $a,b,\ell$ with
    $a+2\ell\le b$ and $b+2\ell\le L-1$.
    Let the label set consist of $L_r,R_r$ for $0\le r<2\ell$, together
    with $B_0,B_1,B_2$. Suppose the supplied operator family $F$ satisfies
    \begin{align}
        \operatorname{supp}F_{L_r} & \subseteq[a+r,a+r+2),         &
        \operatorname{supp}F_{R_r} & \subseteq[b+r,b+r+2),\notag     \\
        \operatorname{supp}F_{B_0} & \subseteq[0,a+1),             &
        \operatorname{supp}F_{B_1} & \subseteq[a+2\ell,b+1),\notag   \\
        \operatorname{supp}F_{B_2} & \subseteq[b+2\ell,L).
        \label{eq:three_region_window_supports}
    \end{align}
    Intervals are intersected with the finite chain. Fix $0\le p,s<2\ell$
    and set
    \begin{align}
        T & =[0,a+p+1),\notag           \\
        E & =[a+p+1,b+s+1),\notag       \\
        R & =V\setminus(T\cup E),\notag \\
        C & =E\setminus[a+2\ell+1,b).
        \label{eq:three_region_selected_regions}
    \end{align}
    There are regional families $H^L,H^M,H^R$ and crossing operators $U,V_0$
    such that, for every label $i\ne L_p,R_s$,
    \begin{align}
        F_i^g                   & =H_i^L\otimes \Id_E\otimes \Id_R
        +\Id_T\otimes H_i^M\otimes \Id_R
        +\Id_T\otimes \Id_E\otimes H_i^R,\notag                      \\
        F_{L_p}^g               & =U\otimes \Id_R,                 &
        F_{R_s}^g               & =\Id_T\otimes V_0,\notag           \\
        \operatorname{OSR}(U)   & \le q^2,                         &
        \operatorname{OSR}(V_0) & \le q^2.
        \label{eq:three_region_two_window_factors}
    \end{align}
    The same choices also satisfy
    \begin{align}
        H^M_{B_0}                    & =H^M_{B_2}=0,                   &
        H^L_{B_1}                    & =H^R_{B_1}=0,\notag               \\
        H^M_{L_r}                    & =0\quad(r\le p),                &
        H^M_{R_r}                    & =0\quad(s\le r),\notag            \\
        \operatorname{supp}H^M_{L_r} & \subseteq C,                    &
        \operatorname{supp}H^M_{R_r} & \subseteq C\quad(0\le r<2\ell).
        \label{eq:three_region_collar_conditions}
    \end{align}
    The three bulk operators are supplied independently; no identity with
    complementary sums of bond terms is assumed.
\end{theorem}
\begin{proof}\leanok
    \uses{thm:chain_two_cut_edge_supports,lem:peps_frame_split_forms,
        thm:peps_supported_left_factor_rank,thm:peps_supported_right_factor_rank,
        thm:three_region_partial_support}
    The selected bonds meet each side of their cut at one site, so their
    crossing operators have operator Schmidt rank at most $q^2$.
    Every other left-window bond lies either in $T$ or in $E$, according as
    its position is before or after $p$. Similarly, every other right-window
    bond lies in $E$ or $R$ according as its position is before or after $s$.
    Choose zero for every unused regional component, including the two selected
    middle components. The three bulk supports lie in $T,E,R$, respectively.
    This gives the identities and zero conditions.

    Each nonzero window middle component inherits its original two-site support
    restricted to $E$, by the partial-matrix-element argument above applied to
    an identity lift. The bounds on the window positions place this restricted
    support in $C$. Monotonicity of support gives the last two assertions.
\end{proof}

\begin{theorem}[The weighted middle sum and its supported remainder]
    \label{thm:three_region_weighted_middle_remainder}
    \lean{TNLean.MPS.AreaLaw.exists_weighted_ordinary_middle_eq_central_add_supported}
    \leanok
    Let $E$ be any finite set, let $H^M$ be indexed by the labels above, and
    let $C\subseteq E$. Suppose $H^M_{B_0}=H^M_{B_2}=0$,
    $H^M_{L_r}=0$ for $r\le p$, and $H^M_{R_r}=0$ for $s\le r$.
    Suppose also that $H^M_{L_r}$ is supported on $C$ for $p<r$, and that
    $H^M_{R_r}$ is supported on $C$ for $r<s$.
    For arbitrary complex weights $z_i$, put
    $O=\{i:i\ne L_p,R_s\}$. There is an operator $K$ supported on $C$ with
    \begin{align}
        \sum_{i\in O}z_iH_i^M=z_{B_1}H^M_{B_1}+K.
        \label{eq:three_region_weighted_middle_remainder}
    \end{align}
\end{theorem}
\begin{proof}\leanok
    Take
    \begin{align}
        K=\sum_{p<r<2\ell}z_{L_r}H^M_{L_r}
        +\sum_{0\le r<s}z_{R_r}H^M_{R_r}.
        \label{eq:three_region_weighted_remainder_formula}
    \end{align}
    The discarded terms vanish by the stated hypotheses. Support on $C$ is
    preserved by scalar multiplication and finite addition.
\end{proof}
